% !TEX root = ../main.tex
\chapter*{Requests about these notes}\label{ch:requests}
\addcontentsline{toc}{chapter}{Requests about these notes}
\markboth{Requests about these notes}{}
\begin{paracol}{2}

Your margin notes are of two kinds: notes on the \emph{content} (a doubt, an idea about a mechanism),
which stay next to the text they refer to, and requests about the \emph{notes themselves} (layout,
boxes, conventions), which do not depend on where they were written. The second kind is collected here,
with what was done about it, so that the chapters read more smoothly. Write new ones anywhere with
\verb|\note{...}| and they will be moved here, or add them directly at the end of this list.

\paragraph{How the margin notes are marked.}
\begin{itemize}
  \item \verb|\note{...}|: a \emph{new} note, not yet read. ``Read my notes'' means these, and only these.
  \item A note that has been answered becomes one box with the question on top and a dark yellow
        \textbf{ANSWER} bar below it (\verb|\begin{answer}{question}[topic]| \dots\ \verb|\end{answer}|). A note with
        several questions gets one ANSWER bar per question. A question asked a second time gets a short
        answer that links to the first one.
  \item \verb|\note*{...}|: a request that has been carried out (green DONE tag).
  \item \verb|\note+{...}|: a comment that has been read and needs no answer (blue SEEN tag).
  \item \verb|\feedback{...}|: a comment on one of your own solutions (indigo).
\end{itemize}

\section*{The list}

\request{done}{All chapters, exercises}{esercizi brevi/a sé stanti lungo il testo, quelli collegati o complessi alla fine}{Standalone exercises now follow the section they practise (right before the next heading); connected or integrative exercises stay together at the end of the chapter, in the original order. In chapters where none were left, the final Exercises section was removed.}
